\section{The Stopped Preference-Formation Economy}
\subsection{Admissibility and the Boundary Contract}
The interior state is $(u,X)\in(1.2,2.8)\times(0.5,2)$, with the controls and coefficients stated in the main paper. The process is stopped at the first exit or $T=1$, whichever occurs first. It receives $G(u,X)=-0.02(u-2)^2+0.1\log X$ at that time. A start on the boundary is settled immediately. This is a finite-horizon contract on an open operating domain; it is not a claim that the unstopped diffusion remains in a rectangle.

For bounded progressively measurable controls, the wealth SDE with coefficients linear in wealth and the additive preference SDE have a well-defined solution up to exit. Before exit, wealth and consumption are bounded away from zero. The flow, terminal payoff, and quadratic adjustment cost are bounded on the stopped domain. Thus the additive objective is finite without introducing a reflected local time or a wealth-transfer process.

The original state-constraint interpretation cannot hold with these coefficients. Preference volatility is $0.05$ at both endpoints, and the maximal wealth drift at the lower boundary is $(0.02+0.8\times0.06)0.5-0.05=-0.016$. At zero portfolio volatility it is more negative. Neither fact is repaired by projecting an Euler endpoint back onto the rectangle. The stopped contract retains all interior coefficients and makes the change in boundary economics explicit.

\subsection{A Boundary-Compatible Payoff}
The choice of $G$ has a useful analytic check. Its derivatives are $G_u=-0.04(u-2)$, $G_{uu}=-0.04$, $G_X=0.1/X$, $G_{XX}=-0.1/X^2$, and $G_{uX}=0$. For every admissible action and every $k\geq0$,
\begin{align*}
 \ell^a+\mathcal L^aG-\rho G
 &\leq \max_{u\in[1.2,2.8]}\frac{0.8^{1-u}}{1-u}
       +0.0064+0.0043-0.00005+0.0032846\\
 &<-0.8162.
\end{align*}
The consumption bound follows because utility is increasing in consumption. Maximization over $u$ at $c=0.8$ gives the upper endpoint; it can also be checked by the explicit derivative. The $0.0043$ term bounds $0.1[0.02+0.06\pi-c/X]$ by using $\pi\leq0.8$, $c/X\geq0.025$. The omitted portfolio diffusion term is nonpositive, and so are adjustment costs. The last term bounds $-\rho G$ uniformly.

Stopped It\^o integration therefore implies $J^p\leq G$ for every admissible policy. To obtain boundary continuity from below, fix for example a constant admissible nonzero portfolio share. Its two-dimensional covariance is uniformly positive definite on the closed rectangle because $X\geq0.5$ and $|\varrho|<1$. The boundary is regular for this fixed diffusion, and its expected exit time tends to zero as the start tends to a boundary point. Bounded rewards and continuity of $G$ imply that this policy's stopped payoff tends to $G$ there. Since its payoff is bounded above by the optimal value and the latter is bounded above by $G$, the value has the declared boundary limit. This argument establishes boundary compatibility; it does not assert global classical differentiability at corners.

\subsection{Transition, Covariance, and First Exit}
For step $h$, form four equiprobable sign pairs $(z_1,z_2)\in\{-1,1\}^2$. Let
\begin{align*}
 d_u&=\theta h+\sigma_u\sqrt h\,z_1,\\
 d_X&=[(r+\pi(\mu-r))X-c]h
      +\pi X\sigma_X\sqrt h\,[\varrho z_1+\sqrt{1-\varrho^2}z_2].
\end{align*}
For an interior state $s$, $\alpha$ is the first fraction in $[0,1]$ at which the line $s+\alpha d$ meets a face, or one if it stays inside. An exit branch has payoff
\[
 \frac{1-e^{-\rho\alpha h}}{\rho}\ell(s,a)
       +e^{-\rho\alpha h}G(s+\alpha d).
\]
A surviving branch instead uses the continuation value at $s+d$ and the full-step flow annuity. At a boundary start the backup is exactly $G(s)$. In particular there is no nonzero continuation from a boundary node. The reference evaluator uses nonnegative bilinear interpolation weights that sum to one on surviving branches.

These signs reproduce conditional first moments and the covariance matrix of the written diffusion. Polynomial tests on $1,u,X,u^2,X^2,uX$ distinguish covariance terms from drift-product terms of order $h$. At the test state and action the covariance is
\[
 \begin{pmatrix}0.0025&-0.002\\-0.002&0.0256\end{pmatrix}.
\]
The largest generator discrepancy decreases by a factor of ten when $h$ decreases from $10^{-2}$ to $10^{-3}$ and again to $10^{-4}$. Preference controls of $-0.2$ and $0.2$ give different mean next states even when both displacements are smaller than a spatial grid cell. No nearest-neighbor map erases the control.

The line-exit rule is an approximation to diffusion killing, not exact Brownian bridge detection. Before-correction overshoots, branch-exit frequencies, occupied exit probabilities, and post-exit payoffs are distinct recorded objects. A reflection-regulator field is inapplicable and is represented by null with that explanation. For the joint state--time refinement, the squared largest mesh width divided by $h$ decreases as $0.04,0.02,0.01$; the separate control refinements are also retained. These observations support consistency diagnostics but do not supply a numerical bound for every $\kappa_n$ in the main theorem.

\subsection{Economic Comparative Statics}
For each admissible policy let $B(p)$ denote cost-free benefit. Define discounted cumulative effort by $C(p)=\E\int_0^{\tau_p}e^{-\rho t}\theta_t^2/2\,dt$. The stopping time is policy dependent but not directly dependent on $k$ for a fixed policy. Thus $V(k)=\sup_p[B(p)-kC(p)]$ is decreasing and convex. For optimizers $p_1,p_2$ at $k_1<k_2$, revealed optimality gives
\begin{align*}
 B(p_1)-k_1C(p_1)&\geq B(p_2)-k_1C(p_2),\\
 B(p_2)-k_2C(p_2)&\geq B(p_1)-k_2C(p_1).
\end{align*}
Adding implies $(k_2-k_1)[C(p_2)-C(p_1)]\leq0$. Also $-C(p)$ is a subgradient of $V$ at any $k$ for which $p$ is optimal, and $V'(k)=-C(p)$ wherever $V$ is differentiable. For policies with value losses at most $e_1,e_2$, the same inequalities with these errors imply
\[
 C(p_2)-C(p_1)\leq\frac{e_1+e_2}{k_2-k_1}.
\]
The result applies to the continuous stopped model and separately to each exactly specified positive finite approximation. A finite-grid comparison cannot set $e_i=0$ for the continuous problem merely because its discrete maximization was exhaustive.

The reported effort is computed both backward and by forwarding the induced state probabilities. The latter multiplies live branch mass by bilinear weights and accumulates discounted cost through the line-exit fraction. This is an occupation calculation for the interpolation-defined finite Markov chain. It provides an economically interpretable sum rather than an unweighted mean of control magnitudes. A reduction of this aggregate does not imply a pointwise reduction of $\theta$, consumption, or portfolio share. Nor does a comparative static identify $k$ from observations without a specified data model and an injective moment map.
